\section{K-012: an antipodal-spectrum criterion for the full stabilizer}

The following finite-group argument resolves the residual determinant
kernel for thirteen of the seventy-seven five-dimensional spectral types.
It also gives a criterion in every odd dimension.  In this section a
tensor is \emph{concise} if each of its three contraction maps is injective.
Positive definite reduced matrices, as in the critical representatives
of K-010, imply conciseness.

\begin{theorem}[K-012]
Let $n$ be odd, and let $T\in(\mathbb C^n)^{\otimes3}$ be a concise
tensor whose product-linear stabilizer $L$ has identity component equal
to the scalar kernel
\[
 S=\{(aI,bI,cI):abc=1\}.
\]
Suppose $g\in L$ has order $M=2^m$, $m\ge1$, and its determinant-product character
$\chi(g)$ has order $M$.  Suppose, moreover, that
\[
                   C_L(g)=S\langle g\rangle.                 \tag{K-012a}
\]
If the spectrum of at least one factor of $g$ contains no pair
$\{\lambda,-\lambda\}$, then
\[
                    L=S\langle g\rangle,
\]
and the effective $SL_n^3$ stabilizer of $T$ is trivial.
\end{theorem}

\begin{proof}
Put $H=L/S$, a finite group.  First
$C_H(\bar g)=\langle\bar g\rangle$, of order $M$.
Indeed, if $hgh^{-1}=sg$ with $s=(s_1I,s_2I,s_3I)\in S$, then
$s_i^M=1$ by taking $M$th powers, and $s_i^n=1$ by taking determinants
in each factor.  Since $(M,n)=1$, all $s_i=1$, and (K-012a) applies.
The character $\chi$ descends to $H$ and is faithful on
$\langle\bar g\rangle$.

Character column orthogonality (Webb, \emph{A Course in Finite Group
Representation Theory}, 23 February 2016 author manuscript,
Corollary~3.4.4, printed p.~36) gives
\[
             \sum_{\rho\in\operatorname{Irr}(H)}
                    |\rho(\bar g)|^2=M.
\]
Every linear character contributes one.  There are at least $M$ linear
characters, namely the first $M$ powers of $\chi$, and hence exactly $M$.
Thus $H^{\rm ab}$ is cyclic of order $M$ and $\chi$ identifies it with
$\mu_M$.  In particular $K=\ker\chi$ is the commutator subgroup of $H$.
No assumption about the order or solvability of $K$ is needed below.

To treat the projective actions of $H$ on the three tensor factors,
take the finite central extension
\[
 \widetilde H=\{(A,B,C)\in L:\det A=\det B=1\}
       \longrightarrow H.
\]
It is surjective, since scalar factors can normalize the first two
determinants, and its kernel is $Z\simeq\mu_n^2$.  It is finite because
each fiber over the finite group $H$ has $n^2$ points.  By multiplying
$g$ by an element of $S$ whose scalar entries lie in $\mu_M$, we may
and do arrange $g\in\widetilde H$ while retaining its order $M$.
This normalization is possible because raising to the $n$th power
permutes $\mu_M$.  It does not change whether a factor spectrum has
an antipodal pair.  We still have
$C_{\widetilde H}(g)=Z\langle g\rangle$, of order $n^2M$.

For a character $\omega\in\widehat Z$, let
\[
 S_\omega(u)=
 \sum_{\substack{\rho\in\operatorname{Irr}(\widetilde H)\\
                  \rho|_Z=(\dim\rho)\omega}}
                         |\rho(g^u)|^2,
 \qquad u\text{ odd}.
\]
For any nonidentity $z\in Z$, the elements $g^u$ and $g^uz$ have
different orders: the first has order $M$, and the second has order
$M\operatorname{ord}(z)$, with $\operatorname{ord}(z)>1$ odd.
They are not conjugate.  Column orthogonality applied to these two
elements, followed by Fourier inversion on $Z$, therefore yields
\[
                         S_\omega(u)=M                     \tag{K-012b}
\]
for every $\omega$ and every odd $u$.  Explicitly the orthogonality
identities are
$\sum_\omega\overline{\omega(z)}S_\omega(u)
=n^2M\,\mathbf1_{z=1}$, and division by $|Z|=n^2$ gives (K-012b).

Fix an irreducible constituent $V$ of one tensor factor, regarded as a
$\widetilde H$-module.  If $V\simeq V\otimes\chi^j$ for some
$0<j<M$, its $g$-spectrum is invariant under multiplication by
$\chi(g)^j$.  Every orbit under this nontrivial two-power root contains
an antipodal pair.  Thus a constituent of an antipodal-free factor has
no such self-twist.  Its $M$ twists by $\chi^j$ are distinct irreducibles
with the same central character.  Equation (K-012b) implies
\[
                         |\operatorname{tr}(g^u|V)|\le1
                         \quad(u\text{ odd}).               \tag{K-012c}
\]

Choose a primitive $M$th root $\zeta$ and let $a_j$ be the multiplicity
of $\zeta^j$ on $V$.  Put $c_j=a_j-a_{j+M/2}$ for $0\le j<M/2$.
Orthogonality of the odd powers of $\zeta$ gives the exact Parseval
identity
\[
 \frac{2}{M}\sum_{\substack{0\le u<M\\u\ {
m odd}}}
       |\operatorname{tr}(g^u|V)|^2
       =\sum_{j=0}^{M/2-1}c_j^2.                            \tag{K-012d}
\]
Indeed the sum of $\zeta^{u(j-k)}$ over odd $u$ is zero for distinct
$j,k$ in this half-range.  By (K-012c), the right side is at most one.
When the spectrum has no antipodal pair,
$\dim V=\sum_j|c_j|$.  The $c_j$ are integers, so $\dim V=1$.
Every irreducible constituent of an antipodal-free factor is therefore
a line.

The center $Z$ acts by the same scalar character on every line of a
fixed tensor factor.  Ratios of their characters consequently descend
to linear characters of $H$, and all such characters are powers of
$\chi$.  Hence $K$ acts projectively trivially on that factor.

Assume this is the first factor.  The subgroup of $L$ fixing that
factor pointwise is connected.  To see this without assuming an
invertible slice, write the first-factor slices of $T$ as matrices
$M_1,\ldots,M_n$.  The finite-dimensional vector space
\[
 E=\{(B,D)\in\operatorname{End}(\mathbb C^n)^2:
                       BM_i=M_iD\text{ for every }i\}
\]
is a unital associative algebra.  Its unit group is exactly the
pointwise stabilizer under $(B,D)\mapsto(B,D^{-t})$ in the last two
factors.  This unit group is a nonempty Zariski-open subset of the
affine vector space $E$, hence connected.  The hypothesis $L^0=S$
forces it to be the one-dimensional scalar subgroup
$(bI,b^{-1}I)$.  Each element of $K$ has a representative in $L$
whose first factor is scalar; multiplying by an element of $S$ makes
that first factor the identity.  The preceding connectedness shows
that the representative then lies in $S$, so $K$ is trivial.
Finally $H\simeq\mu_M$ is generated by $\bar g$.
An $SL_n^3$ stabilizer element has $\chi=1$, hence lies in $S$;
its class in the effective group is therefore trivial.  This proves
the theorem.
\end{proof}

\subsection*{Application to K-010}
For the simple-spectrum types, (K-012a) follows from diagonal
centralization and Smith form $(1^{12},16)$.  For the repeated-spectrum
type, a centralizing stabilizer element preserves each of the three
distinct projective lines in its two-dimensional eigenspace, so its
restriction to that eigenspace is scalar; the same diagonal Smith
calculation applies.  The exact rank-$73$ calculation in K-010 gives
$L^0=S$ for every type.  Each tensor is concise because it has a critical
representative with identity reduced matrices.

Exactly thirteen of the seventy-seven spectral types have an
antipodal-free factor.  Their indices in the K-010 table are
\[
        1,25,38,39,40,41,42,43,50,51,67,68,73.                \tag{K-012e}
\]
For these types the full stabilizer is now proved to be
$S\times\mu_{16}$, with faithful determinant character; the effective
$SL_5^3$ stabilizer is trivial.  Their quotient points are smooth and
their fixed tangent dimensions are one.  The fixed tensor dimensions
are $13$ and the fixed full Lie dimensions are $15$ in the simple cases;
for the repeated case they are $15$ and $17$.  Subtracting the
determinant-one Lie dimension in either case gives one.

The remaining sixty-four types have an antipodal pair in each factor.
This theorem does not decide their full finite stabilizers.  The chain
type, number $59$, is covered separately by the rechecked P-F60.
